% GENERATED FILE. Edit canonical lessons, then run npm run generate:books.
% canonical-source-hash: fnv1a64:3fc26096a601e184
% canonical-lessons: MR-C430-mobadla, MR-C430-dengi, MR-C430-vargani, MR-C430-hissa, MR-C430-urne

\chapter{Payments and Donations}
\label{ch:mr-payments-and-donations}
\hlchaptermodality{430}

\begin{chapteropening}
\textbf{By the end of this chapter:} \emph{I can talk about payment for work, donations, contributions and shares, and what is left over.}

Complete the last lesson of chapter 430: I can talk about payment for work, donations, contributions and shares, and what is left over.
\end{chapteropening}

\section[\texorpdfstring{\mr{मोबदला} (mobadlā)}{मोबदला (mobadlā)}]{\mr{मोबदला} (mobadlā) — a payment in return}

\label{lesson:MR-C430-mobadla}



\begin{quote}
Before the new one: say the Marathi for to recover money, then the Marathi for a lease.
\end{quote}

\subsection*{You'll want to know: \mr{मोबदला}}
\textbf{\mr{मोबदला}} — \emph{mobadlā} — \textquotedblleft{}a payment in return\textquotedblright{}.

\begin{quote}
\mr{कामाचा} \mr{योग्य} \mr{मोबदला} \mr{मिळाला} \mr{पाहिजे}.
\end{quote}

\begin{quote}
\emph{kāmācā yogya mobadlā miḷālā pāhije.}
\end{quote}

\begin{quote}
Work should get a fair payment.
\end{quote}

Masculine.

A fee or compensation for work or for land taken.

\begin{cousinweb}
\textbf{\mr{मोबदला}} (\emph{mobadlā}) is a loan through Persian, usually traced to Arabic \emph{mubādala}, \textquotedblleft{}an exchange\textquotedblright{}.
\end{cousinweb}

\begin{grammarlens}[title={What you've built}]
One more word for getting around.
\end{grammarlens}

\subsection*{Your turn}
\begin{itemize}
\raggedright
  \item \emph{Say it:} \emph{mobadlā}
  \item \emph{Say it:} \emph{mobadlā}, once more
  \item \emph{Say it:} \emph{kāmācā yogya mobadlā miḷālā pāhije.} — Work should get a fair payment.
  \item \emph{Recall:} say the Marathi for to recover money, then the Marathi for a lease, then say \emph{mobadlā} again
\end{itemize}

\begin{tcolorbox}[breakable,colback=teal!4,colframe=teal!35!black,title={Before you move on}]
What does \mr{मोबदला} mean? (\textquotedblleft{}A payment in return\textquotedblright{}.) Say it once more.
\end{tcolorbox}
\section[\texorpdfstring{\mr{देणगी} (deṇgī)}{देणगी (deṇgī)}]{\mr{देणगी} (deṇgī) — a donation}

\label{lesson:MR-C430-dengi}



\begin{quote}
Before the new one: say the Marathi for a lease, then the Marathi for a payment in return.
\end{quote}

\subsection*{You'll want to know: \mr{देणगी}}
\textbf{\mr{देणगी}} — \emph{deṇgī} — \textquotedblleft{}a donation\textquotedblright{}.

\begin{quote}
\mr{त्यांनी} \mr{शाळेला} \mr{मोठी} \mr{देणगी} \mr{दिली}.
\end{quote}

\begin{quote}
\emph{tyānnī śāḷelā moṭhī deṇgī dilī.}
\end{quote}

\begin{quote}
They gave the school a large donation.
\end{quote}

Feminine.

Built on \mr{देणे}, to give; the receiver takes -\mr{ला}.

\begin{cousinweb}
\textbf{\mr{देणगी}} (\emph{deṇgī}) is built on \textbf{\mr{देणे}}, \textquotedblleft{}to give\textquotedblright{}, with a noun ending; the shape of the ending is not settled.
\end{cousinweb}

\begin{grammarlens}[title={What you've built}]
One more word for getting around.
\end{grammarlens}

\subsection*{Your turn}
\begin{itemize}
\raggedright
  \item \emph{Say it:} \emph{deṇgī}
  \item \emph{Say it:} \emph{deṇgī}, once more
  \item \emph{Say it:} \emph{tyānnī śāḷelā moṭhī deṇgī dilī.} — They gave the school a large donation.
  \item \emph{Recall:} say the Marathi for a lease, then the Marathi for a payment in return, then say \emph{deṇgī} again
\end{itemize}

\begin{tcolorbox}[breakable,colback=teal!4,colframe=teal!35!black,title={Before you move on}]
What does \mr{देणगी} mean? (\textquotedblleft{}A donation\textquotedblright{}.) Say it once more.
\end{tcolorbox}
\section[\texorpdfstring{\mr{वर्गणी} (vargaṇī)}{वर्गणी (vargaṇī)}]{\mr{वर्गणी} (vargaṇī) — a subscription; a contribution}

\label{lesson:MR-C430-vargani}



\begin{quote}
Before the new one: say the Marathi for a payment in return, then the Marathi for a donation.
\end{quote}

\subsection*{You'll want to know: \mr{वर्गणी}}
\textbf{\mr{वर्गणी}} — \emph{vargaṇī} — \textquotedblleft{}a subscription; a contribution\textquotedblright{}.

\begin{quote}
\mr{उत्सवासाठी} \mr{सगळ्यांनी} \mr{वर्गणी} \mr{दिली}.
\end{quote}

\begin{quote}
\emph{utsavāsāṭhī sagaḷyānnī vargaṇī dilī.}
\end{quote}

\begin{quote}
Everyone gave a contribution for the festival.
\end{quote}

Feminine.

A share each member pays: \mr{मासिक} \mr{वर्गणी}, a monthly subscription.

\begin{cousinweb}
\textbf{\mr{वर्गणी}} (\emph{vargaṇī}) is a Marathi word for a share paid into a common fund; its origin is not settled.
\end{cousinweb}

\begin{grammarlens}[title={What you've built}]
One more word for getting around.
\end{grammarlens}

\subsection*{Your turn}
\begin{itemize}
\raggedright
  \item \emph{Say it:} \emph{vargaṇī}
  \item \emph{Say it:} \emph{vargaṇī}, once more
  \item \emph{Say it:} \emph{utsavāsāṭhī sagaḷyānnī vargaṇī dilī.} — Everyone gave a contribution for the festival.
  \item \emph{Recall:} say the Marathi for a payment in return, then the Marathi for a donation, then say \emph{vargaṇī} again
\end{itemize}

\begin{tcolorbox}[breakable,colback=teal!4,colframe=teal!35!black,title={Before you move on}]
What does \mr{वर्गणी} mean? (\textquotedblleft{}A subscription; a contribution\textquotedblright{}.) Say it once more.
\end{tcolorbox}
\section[\texorpdfstring{\mr{हिस्सा} (hissā)}{हिस्सा (hissā)}]{\mr{हिस्सा} (hissā) — a share; a portion}

\label{lesson:MR-C430-hissa}



\begin{quote}
Before the new one: say the Marathi for a donation, then the Marathi for a subscription.
\end{quote}

\subsection*{You'll want to know: \mr{हिस्सा}}
\textbf{\mr{हिस्सा}} — \emph{hissā} — \textquotedblleft{}a share; a portion\textquotedblright{}.

\begin{quote}
\mr{तिघा} \mr{भावांनी} \mr{शेताचे} \mr{तीन} \mr{हिस्से} \mr{केले}.
\end{quote}

\begin{quote}
\emph{tighā bhāvānnī śetāce tīn hisse kele.}
\end{quote}

\begin{quote}
The three brothers split the field into three shares.
\end{quote}

Masculine. Plural: \mr{हिस्से}, \emph{hisse}.

The plural is \mr{हिस्से}.

\begin{cousinweb}
\textbf{\mr{हिस्सा}} (\emph{hissā}) is a loan through Persian, usually traced to Arabic \emph{ḥiṣṣa}, \textquotedblleft{}a portion, a share\textquotedblright{}.
\end{cousinweb}

\begin{grammarlens}[title={What you've built}]
One more word for getting around.
\end{grammarlens}

\subsection*{Your turn}
\begin{itemize}
\raggedright
  \item \emph{Say it:} \emph{hissā}
  \item \emph{Say it:} \emph{hissā}, once more
  \item \emph{Say it:} \emph{tighā bhāvānnī śetāce tīn hisse kele.} — The three brothers split the field into three shares.
  \item \emph{Recall:} say the Marathi for a donation, then the Marathi for a subscription, then say \emph{hissā} again
\end{itemize}

\begin{tcolorbox}[breakable,colback=teal!4,colframe=teal!35!black,title={Before you move on}]
What does \mr{हिस्सा} mean? (\textquotedblleft{}A share; a portion\textquotedblright{}.) Say it once more.
\end{tcolorbox}
\section[\texorpdfstring{\mr{उरणे} (urṇe)}{उरणे (urṇe)}]{\mr{उरणे} (urṇe) — to be left over}

\label{lesson:MR-C430-urne}



\begin{quote}
Before the new one: say the Marathi for a subscription, then the Marathi for a share.
\end{quote}

\subsection*{You'll want to know: \mr{उरणे}}
\textbf{\mr{उरणे}} — \emph{urṇe} — \textquotedblleft{}to be left over\textquotedblright{}.

\begin{quote}
\mr{जेवणानंतर} \mr{थोडा} \mr{भात} \mr{उरला}.
\end{quote}

\begin{quote}
\emph{jevṇānantar thoḍā bhāt urlā.}
\end{quote}

\begin{quote}
After the meal a little rice was left over.
\end{quote}

Past (he …): \mr{तो} \mr{उरला}, \emph{to urlā}.

What is left is the subject: \mr{पैसे} \mr{उरले}, money was left.

\begin{cousinweb}
\textbf{\mr{उरणे}} (\emph{urṇe}) is inherited: CDIAL lists Marathi \emph{urṇẽ}, \textquotedblleft{}to be left as a residue\textquotedblright{}, under \emph{*urvara}, \textquotedblleft{}surplus, left over\textquotedblright{} (CDIAL 2356).
\end{cousinweb}

\begin{grammarlens}[title={What you've built}]
One more word for getting around.
\end{grammarlens}

\subsection*{Your turn}
\begin{itemize}
\raggedright
  \item \emph{Say it:} \emph{urṇe}
  \item \emph{Say it:} \emph{urṇe}, once more
  \item \emph{Say it:} \emph{jevṇānantar thoḍā bhāt urlā.} — After the meal a little rice was left over.
  \item \emph{Recall:} say the Marathi for a subscription, then the Marathi for a share, then say \emph{urṇe} again
\end{itemize}

\begin{tcolorbox}[breakable,colback=teal!4,colframe=teal!35!black,title={Before you move on}]
What does \mr{उरणे} mean? (\textquotedblleft{}To be left over\textquotedblright{}.) Say it once more.
\end{tcolorbox}
